\section{Appendix: Polynomial construction and the finite-sample certificate}

The achievability argument for \cref{obj:thm:finite-certificate} combines endpoint concentration, control of polynomial derivatives, and exact Gaussian second moments. These ingredients connect the observable correction in \cref{obj:def:inverse-heat} to the stabilized ratios and public intervals in \cref{obj:def:ratio,obj:synth_6}. We present the auxiliary results in that mathematical order.

\subsection{Hermite polynomials and Gaussian averaging}

The Gaussian identities use the probabilists' Hermite polynomials. We fix their normalization before stating the orthogonality result.



This normalization determines the factorial weights in Gaussian polynomial second moments. The relevant classical identity is the following.

\begin{theoremv}{lem:classical-hermite-facts}[Gaussian Hermite orthogonality]
Let \(H_j\) be the probabilists' Hermite polynomials, defined by
\[
\exp(zt-t^2/2)
=\sum_{j=0}^{\infty}H_j(z)\frac{t^j}{j!}.
\]
For a standard Gaussian variable \(Z\) and nonnegative integers \(j,k\),
\[
\mathbb E[H_j(Z)H_k(Z)]
=j!\mathbf1\{j=k\},
\qquad H_0=1.
\]
\end{theoremv}

Orthogonality separates the contributions of distinct polynomial orders under Gaussian averaging. Consequently, a finite Hermite expansion has a second moment expressed as a sum of squared coefficients with factorial weights. This structure supplies the exact second-moment identity in \cref{obj:lem:exact-covariance}.

\subsection{Positive endpoint concentration}

The latent averaging step uses the endpoint polynomial weight in \cref{obj:def:endpoint-kernel}. Its positivity and concentration determine how accurately a weighted arm mean approximates the corresponding cutoff mean. The following bounds quantify that concentration and its squared-integral cost.

\begin{lemmav}{lem:positive-endpoint}[Positive endpoint kernel bounds]
Let \(0<\beta\le1\), and let \(K_L\) be the endpoint kernel defined in \cref{obj:def:endpoint-kernel}. Define its bias moment by
\[
M_{L,\beta}=\int_0^1 x^\beta K_L(x)\,dx.
\]
There exists a constant \(C>0\), depending on \(\beta\), such that for every integer \(L\ge1\),
\[
K_L(x)\ge0\quad\text{for all }x\in[0,1],
\qquad
\int_0^1 K_L(x)\,dx=1,
\]
and
\[
M_{L,\beta}\le C L^{-2\beta},
\qquad
\int_0^1 K_L(x)^2\,dx\le C L^2.
\]
\end{lemmav}

Positivity and unit integral make the latent weight a normalized averaging measure on the reflected arm interval. Its Hölder moment decreases at order \(L^{-2\beta}\), giving the endpoint construction spatial resolution of order \(L^{-2}\). The squared-integral bound records the corresponding sampling cost before measurement-error correction.

The density and mean restrictions enter distinct parts of this averaging argument. The density lower bound in \cref{obj:ass:density-bounds} supplies the population denominator bound \(B_{d,L}(P)\ge1/4\) in \cref{obj:thm:finite-certificate}. The upper density bound controls the amount of mass multiplying the polynomial weight. Variation of the potential-outcome mean around its cutoff value is controlled by \cref{obj:ass:mean-holder}, with \(M_{L,\beta}\) measuring the resulting endpoint approximation error. Thus the approximation term is governed by potential-mean smoothness, while the density bounds control normalization and weighted mass.

\subsection{Derivative growth and inverse-heat second moments}

Gaussian correction introduces derivatives of the endpoint polynomial. Their growth determines how measurement error changes the second-moment cost. The required inequality applies to a real polynomial \leanref{sym:Polynomial}{\ensuremath{p}} of bounded degree.

\begin{lemmav}{lem:factorial-derivatives}[Factorial polynomial derivative bound]
Let \(m,j\) be nonnegative integers with \(1\le j\le m\), and let \(p\) be a real polynomial of degree at most \(m\). Writing \(p^{(j)}\) for its \(j\)-th ordinary derivative, we have
\[
\left(\int_0^1 \bigl(p^{(j)}(x)\bigr)^2\,dx\right)^{1/2}
\le
\frac{4^j(m+1)^{2j}}{j!}
\left(\int_0^1 p(x)^2\,dx\right)^{1/2}.
\]
\end{lemmav}

The factorial dependence on \leanref{sym:j}{\ensuremath{j}}, the derivative order, controls the accumulation of high-order terms in the inverse-heat second moment. Together with the squared-integral bound in \cref{obj:lem:positive-endpoint}, it gives a degree-dependent bound for every derivative contributing to the correction.

The inverse-heat polynomial has both an exact conditional expectation and an exact conditional second moment. These identities apply at every real latent-score argument and throughout the declared noise range.

\begin{lemmav}{lem:exact-covariance}[Exact Gaussian moment identities]
There exists an absolute constant \(C>0\) such that, for every integer \(L\ge1\) and every \(\sigma\in[0,1]\), the following holds. Let \(K_L\) and \(Q_{L,\sigma}\) be defined by \cref{obj:def:endpoint-kernel,obj:def:inverse-heat}, and set
\[
m_L=3(L-1),\qquad
V_L(\sigma)=\frac34\sum_{j=0}^{m_L}\frac{\sigma^{2j}}{j!}
\int_0^1 [K_L^{(j)}(x)]^2\,dx.
\]
For a standard Gaussian random variable \(\varepsilon\), at every \(x\in\mathbb R\),
\[
\mathbb E\!\left[Q_{L,\sigma}(x+\sigma\varepsilon)\right]=K_L(x),
\qquad
\mathbb E\!\left[Q_{L,\sigma}(x+\sigma\varepsilon)^2\right]
=\sum_{j=0}^{m_L}\frac{\sigma^{2j}[K_L^{(j)}(x)]^2}{j!}.
\]
Moreover,
\[
V_L(\sigma)\le C L^2
\sum_{j=0}^{m_L}\frac{(C\sigma^2L^4)^j}{(j!)^3}.
\]
\end{lemmav}

The first identity identifies the latent endpoint average targeted by the observable correction. The second accounts for the full Gaussian distribution of the proxy, including its tails at positive noise scales. The cubic factorial in the final bound combines the factorial derivative control with the Gaussian second-moment weights.

Under \cref{obj:ass:gaussian,obj:ass:error-independence}, these identities apply conditionally on the latent score and potential outcomes, using the reflection convention in \cref{obj:synth_5}. The density upper bound makes \(V_L(\sigma)\) a common second-moment bound for the arm summands. Binary outcomes also place the outcome-marked summands under that bound. Accordingly, the sampling term in \cref{obj:thm:finite-certificate} accounts for both empirical averages over the entire observed experiment.

\subsection{The finite-sample certificate}

The auxiliary results supply the endpoint approximation and sampling controls used by \cref{obj:thm:finite-certificate}. The construction in \cref{obj:def:ratio} addresses the remaining issue of turning empirical marked and unmarked averages into arm estimates. Its population denominator is bounded below by \(1/4\), while the empirical denominator is floored at \(1/8\). The floor therefore lies below the guaranteed population mass and keeps the ratio defined for every observed sample. Projection onto \([1/4,3/4]\) aligns each arm estimate with the maintained mean range in \cref{obj:ass:mean-bounds}.

For every fixed positive index, \cref{obj:thm:finite-certificate} bounds expected absolute effect-estimation error by
\[
12M_{L,\beta}+28\sqrt{\frac{V_L(\sigma)}{n}}.
\]
The public radius in \cref{obj:synth_6} uses the same approximation term and the larger sampling term \(28\sqrt{40V_L(\sigma)/n}\), delivering coverage of at least \(9/10\). The certificate thus translates positive latent averaging and exact observable second moments into finite-sample estimation and inference guarantees under the benchmark restrictions.

Interval geometry and measurability are also part of the declared decision problem. The projected arm estimates place the estimated contrast in \([-1/2,1/2]\). Hence the intersection defining \(I_L\) in \cref{obj:synth_6} is a nonempty connected closed interval. The target belongs to the same contrast range under \cref{obj:ass:mean-bounds}, so intersection with \([-1,1]\) preserves its coverage. Finite polynomial averages, the denominator floor, projection, and the interval endpoints are measurable functions of the observed sample. Their outputs are constant in the independent procedure seed, placing them within the randomized experiment of \cref{obj:def:honest-class}.

Finally, \cref{obj:def:public-selection} compares radii over a finite public candidate set and resolves ties by choosing the smallest index. For fixed public inputs, the selected index is fixed across samples; its coverage is therefore the coverage of its selected candidate. The index-zero values in \cref{obj:synth_6} provide an estimate of zero and the interval \([-1/2,1/2]\), which contains every admissible target. This fallback supplies coverage one, absolute error at most \(1/2\), and interval length one.

Including index zero ensures \(E_\beta(n,\sigma)\le1/2\). For a selected positive index, the candidate risk bound is bounded by its public radius, and interval length is bounded by twice that radius. These guarantees, together with the fallback, are the selected risk and expected-length bounds in \cref{obj:thm:finite-certificate}. The selected interval consequently belongs to the uniformly honest class in \cref{obj:def:honest-class}, with precision controlled by a deterministic certificate over the full benchmark.
